\documentclass[12pt]{extarticle}

\usepackage{tabularx}
\usepackage{amsmath}
\usepackage{extsizes}
\usepackage{graphicx}
\usepackage{amssymb}
\usepackage[english,russian]{babel}
\usepackage[utf8]{inputenc}
\usepackage[a4paper]{geometry}
\geometry{top=20mm,bottom=30mm,left=20mm,right=15mm}
\usepackage[T1,T2A]{fontenc}
\usepackage{setspace}
\usepackage{indentfirst}
\usepackage{caption}
\renewcommand{\baselinestretch}{1.5}
\usepackage{spverbatim}
\usepackage{listings}
\usepackage{xcolor}
\usepackage{tabularray}

\lstset{
    inputencoding=utf8,
    extendedchars=\true,
    basicstyle=\ttfamily\small,
    breaklines=\true,
    numbers=left,
    numberstyle=\tiny\color{gray},
    frame=single,
    captionpos=b,
    keepspaces=\true,
    showstringspaces=\false
}

\definecolor{codebg}{gray}{0.96}
\definecolor{comment}{gray}{0.45}
\definecolor{keyword}{RGB}{0, 75, 150}
\definecolor{string}{RGB}{150, 80, 30}
\definecolor{number}{gray}{0.5}
\definecolor{identifier}{rgb}{0,0,0}

\lstdefinestyle{coolprint}{
  basicstyle   = \ttfamily\footnotesize,
  keywordstyle = \bfseries\color{keyword},
  commentstyle = \itshape\color{comment},
  stringstyle  = \color{string},
  identifierstyle = \color{identifier},
  backgroundcolor = \color{codebg},
  numbers      = left,
  numberstyle  = \footnotesize\color{number},
  numbersep    = 5pt,
  frame        = single,
  framerule    = 0.8pt,
  framesep     = 3pt,
  breaklines   = true,
  breakatwhitespace = true,
  tabsize      = 4,
  showspaces   = false,
  showstringspaces = false,
  showtabs     = false,
  xleftmargin  = 10pt,
  xrightmargin = 5pt,
  aboveskip    = \medskipamount,
  belowskip    = \medskipamount,
  captionpos   = b,
}
\lstset{style=coolprint}

\begin{document}

\begin{center}
  \textbf{МИНОБРНАУКИ РОССИИ} \\
  \textbf{Санкт-Петербургский государственный} \\
  \textbf{электротехнический университет} \\
  \textbf{«ЛЭТИ» им. В.И. Ульянова (Ленина)} \\
  \textbf{Кафедра ВТ} \\
\end{center}

\vspace{0.25\textheight}

\begin{center}
\textbf{ОТЧЕТ} \\
\textbf{по лабораторной работе №4} \\
\textbf{по дисциплине «ОргЭВМиС»} \\
\textbf{Тема: ИССЛЕДОВАНИЕ СТРАТЕГИЙ КЭШИРОВАНИЯ В RIPES} \\
\end{center}

\vspace*{\fill}

\begin{center}
\begin{tabularx}{\textwidth}{ l X r }
 Студенты гр. 5352 & \rule{8cm}{0.2mm} & Соболевский М.А. \\
 & \rule{8cm}{0.2mm} & Решетников С.А. \\
 Преподаватель & \rule{8cm}{0.2mm} & Кочетков А.В. \\
\end{tabularx}
\end{center}

\vspace{3em}

\begin{center}
Санкт-Петербург\\
\the\year
\end{center}

\thispagestyle{empty}

\pagebreak

\section{Задание на лабораторную работу}

Вариант №8.

Вычислить сумму всех нулевых и единичных битов в двоичном представлении элементов массива \texttt{А} из 500 элементов, запомнить эти значения в памяти и вывести эти значения или сообщение об отсутствии некоторых видов
битов.

Проанализировать работу программы в условиях симуляции кэша в Ripes
(вкладка Cache по левой стороне интерфейса) на небольшом размере кэш памяти
данных (256 бит – data bits) в зависимости от различных стратегий. Для каждого
запуска программы в отчете представить значения следующих метрик: \texttt{Hit Rate},
\texttt{Writebacks}, \texttt{Hits}, \texttt{Misses}.

\section{Текст программы}

\renewcommand{\baselinestretch}{1}
\input{assets/source}
\renewcommand{\baselinestretch}{1.5}

\section{Пример запуска программы}

Вывод программы с массивом \texttt{A} заданным в самом исходнике:

\renewcommand{\baselinestretch}{1}
\begin{spverbatim}
zero bits: 7923
one bits: 8077

Program exited with code: 0
\end{spverbatim}
\renewcommand{\baselinestretch}{1.5}

Все последующие запуски программы с разными конфигурациями кэш-памяти имеют идентичный вывод.

\begin{longtblr}{colspec={|X|l|l|l|l|}, cells={font=\linespread{1.0}\selectfont}, hlines, vlines}
  Preset & Hit Rate & Writebacks & Hits & Misses \\
  1-entry fully associative & 0.942 & 1001 & 32501 & 2002 \\
  2-entry fully associative & 0.971 & 499 & 33503 & 1000 \\
  2-entry direct mapped & 0.971 & 499 & 33503 & 1000 \\
  4-entry fully associative & 0.9927 & 0 & 34250 & 253 \\
  4-entry 2-word 2-way set associative & 0.9819 & 249 & 33877 & 626 \\
  4-entry direct mapped & 0.9765 & 373 & 33691 & 812 \\
  8-entry fully associative & 0.9854 & 0 & 33999 & 504 \\
  8-entry 1-word 4-way set associative & 0.9854 & 0 & 33999 & 504 \\
  8-entry 1-word 2-way set associative & 0.9818 & 125 & 33874 & 629 \\
  8-entry direct mapped & 0.9782 & 249 & 33750 & 753
\end{longtblr}

Теперь с \\
Wr. hit: Write-through \\
Wr. miss: No write allocate

\begin{longtblr}{colspec={|X|l|l|l|l|}, cells={font=\linespread{1.0}\selectfont}, hlines, vlines}
  Preset & Hit Rate & Writebacks & Hits & Misses \\
  4-entry fully associative & 0.9926 & 17001 & 34249 & 254 \\
  4-entry 2-word 2-way set associative & 0.9818 & 17001 & 33876 & 627 \\
  4-entry direct mapped & 0.9692 & 17001 & 33440 & 1063
\end{longtblr}

Теперь с \\
Wr. hit: Write-back \\
Wr. miss: Write allocate \\
Repl. policy: Random

\begin{longtblr}{colspec={|X|l|l|l|l|}, cells={font=\linespread{1.0}\selectfont}, hlines, vlines}
  Preset & Hit Rate & Writebacks & Hits & Misses \\
  8-entry fully associative & 0.9735 & 404 & 33590 & 913 \\
  8-entry 1-word 4-way set associative & 0.9751 & 354 & 33645 & 858 \\
  8-entry 1-word 2-way set associative & 0.9745 & 374 & 33624 & 879 \\
  8-entry direct mapped & 0.9782 & 249 & 33750 & 753
\end{longtblr}

\subsection{Выбор стратегии кэширования и обоснование}

При анализе учитывался одинаковый объём кэша данных (\(256\) бит) и разный
размер памяти тегов для рассмотренных конфигураций. Увеличение числа входов
кэша при том же объёме данных приводит к росту числа теговых записей; кроме
того, fully associative требует больше теговых битов на вход из-за отсутствия
индексных битов. Поэтому конфигурации сравнивались не только по Hit Rate, но
и с учётом роста памяти тегов.

\begin{longtblr}{colspec={|l|c|c|c|X|}, cells={font=\linespread{1.0}\selectfont}}
\hline
Конфигурация & Hit Rate & Misses & Writebacks & Комментарий \\
\hline
4-entry fully associative, WB/WA, LRU & 0.9927 & 253 & 0 & лучший Hit Rate \\
4-entry 2-word 2-way SA, WB/WA & 0.9819 & 626 & 249 & компромисс по памяти тегов \\
4-entry direct mapped, WB/WA & 0.9765 & 812 & 373 & меньше тегов, но больше промахов \\
8-entry fully associative, WB/WA & 0.9854 & 504 & 0 & вдвое больше тегов, Hit Rate ниже \\
4-entry fully associative, WT/NA & 0.9926 & 254 & 17001 & много записей в память \\
8-entry fully associative, Random & 0.9735 & 913 & 404 & Random ухудшает Hit Rate \\
\hline
\end{longtblr}

\paragraph{Выбор стратегии.}
Наиболее перспективной по коэффициенту попаданий является стратегия
\textbf{4-entry fully associative с политикой write-back и write-allocate}.
Для ассоциативных конфигураций предпочтительна политика замещения LRU:
Random replacement снижает Hit Rate, например, для 8-entry fully associative
с \(0.9854\) до \(0.9735\).

\paragraph{Обоснование.}
4-entry fully associative обеспечивает максимальный Hit Rate \(0.9927\) и
минимальное число промахов \(253\). Он превосходит 4-entry 2-way set
associative на \(1.08\) п.п. и 4-entry direct mapped на \(1.62\) п.п. При этом
4-entry FA имеет только 4 входа тегов, тогда как 8-entry FA --- 8 входов, но
проигрывает по Hit Rate (\(0.9854\)). Следовательно, рост памяти тегов при
переходе к 8-entry не компенсируется улучшением попаданий. Write-through /
no write allocate для 4-entry FA почти не меняет Hit Rate (\(0.9926\)), но
даёт \(17001\) writebacks против \(0\) у write-back, что резко увеличивает
трафик записи. Поэтому выбранная стратегия:
\textbf{4-entry fully associative; write-back; write-allocate; LRU}.
Если память тегов критична, допустим компромисс --- 4-entry 2-way set
associative, но по критерию Hit Rate он уступает выбранной конфигурации.

\section{Выводы}

В ходе лабораторной работы была разработана программа для вычисления суммы
нулевых и единичных битов в двоичном представлении элементов массива
\texttt{A} из 500 элементов. Для заданного в исходном коде массива программа
корректно выдала значения \texttt{zero bits: 7923} и \texttt{one bits: 8077}.
При всех последующих запусках с различными конфигурациями кэш-памяти вывод
программы оставался идентичным, что подтверждает независимость результата
вычислений от параметров симуляции кэша.

Анализ метрик кэш-памяти проводился при одинаковом объёме кэша данных
(\(256\) бит) и различных стратегиях отображения, записи и замещения. При этом
учитывался разный размер памяти тегов: fully associative требует больше
теговых битов на вход из-за отсутствия индексных битов, а увеличение числа
входов кэша при том же объёме данных также приводит к росту числа теговых
записей.

Наиболее высокий коэффициент попаданий показала конфигурация
\textbf{4-entry fully associative} с политикой \textbf{write-back} и
\textbf{write-allocate}: \(\text{Hit Rate} = 0.9927\), \(253\) промаха и
\(0\) writebacks. Она превзошла 4-entry 2-way set associative
(\(0.9819\)) и 4-entry direct mapped (\(0.9765\)). Увеличение числа входов
до 8 в fully associative не дало выигрыша: Hit Rate снизился до \(0.9854\)
при вдвое большем числе теговых записей. Следовательно, рост памяти тегов
в рассмотренных конфигурациях не компенсируется улучшением попаданий.

Использование политики \textbf{write-through / no write allocate} для
4-entry fully associative почти не изменило Hit Rate (\(0.9926\)), однако
привело к \(17001\) writebacks против \(0\) при write-back, что резко
увеличивает трафик записи в память. Политика замещения \textbf{Random}
для ассоциативных конфигураций ухудшила Hit Rate: например, для 8-entry
fully associative — с \(0.9854\) до \(0.9735\). Для direct mapped политика
замещения не влияет, поскольку в каждом наборе имеется только один вход.

Таким образом, наиболее перспективной по коэффициенту попаданий является
стратегия \textbf{4-entry fully associative; write-back; write-allocate; LRU}.
Она обеспечивает максимальный Hit Rate, минимальное число промахов и не
создаёт дополнительного трафика writebacks. Если память тегов критична,
допустим компромиссный вариант — 4-entry 2-way set associative, однако по
критерию Hit Rate он уступает выбранной конфигурации. Цель лабораторной
работы достигнута: программа реализована и проверена, работа кэш-памяти
проанализирована для различных стратегий, выбрана и обоснована наиболее
эффективная конфигурация.

\end{document}
